\section{The precise obstruction to a finite-character mixed scalar packet}
\label{sec:mixed-packet-gate}

The constructions above realize the coefficient geometrically.
The following comparison concerns the actual automorphic data in
the supplied source, at the pinned commit already specified.
It rules out two proposed substitutions inside that data.  It
does not rule out a theorem comparing two separate regular packets.

\begin{proposition}[AC-G37: good rows and coincident scalar centers]
\label{AC-G37}
The following statements hold.

First, under the source hypothesis
\path{SumRepresentations.At.TieLab}, a parameter
\texttt{Sig} satisfying \texttt{HasParamW} for a good row is
exactly one cuspidal block of multiplicity one:
\[
                  \operatorname{Sig}=[(r,\pi,1)].       \tag{G37.1}
\]
Consequently an additional scalar Speh block cannot be inserted
inside the existing good rank-$r$ parameter.

Second, no parameter satisfying the literal
\texttt{NBC.infW\_alg} and \texttt{IsBCV} conditions can contain
two positive-length scalar Speh blocks having the same
archimedean center at an embedding.  In particular this excludes
two scalar characters whose ratio is a finite-order character,
including $\chi_{269}$, irrespective of the two lengths and
of the ambient rank.

Third, replacing one of two conjugate-self-dual scalar characters
by a nonzero real norm twist cannot preserve the source
conjugate-self-duality condition.  Thus such a norm shift does
not repair the preceding obstruction within the stated data.
Changes to nonparallel algebraic infinity types are outside this
assertion and require a new coefficient comparison.
\end{proposition}
\begin{proof}
\proofstep{1. Degree saturation and the good-row condition.}
Here is the full polynomial argument, including the hypothesis
that supplies the inertia dictionary.  Write the blocks of
\texttt{Sig} as $(a_i,\pi_i,d_i)$, with $a_i,d_i\ge1$ and
$\sum_i a_id_i=r$.  At a test place $v$, let $h_{v,i}$ be the
source polynomial \texttt{items\_2} for $\pi_i$; the subscript
here denotes the source Unicode subscript $2$.
The two clauses of \texttt{At.TieLab} say that real unramified
twists do not change this polynomial and that
$\deg h_{v,i}\le a_i$.
The exact source lemmas
\texttt{inertiaCharpoly\_speh\_eq\_pow} and
\texttt{inertiaCharpoly\_eq\_of\_iso} therefore give
\[
                    f_v=\prod_i h_{v,i}^{d_i}.        \tag{G37.2}
\]
Choose one test prime $g(k)$ different from the coefficient prime,
which is possible because the source has two distinct test primes.
For the rest of this step, $v$ ranges over \texttt{N.PlG k}
for this single fixed $k$.  At this prime
\texttt{plGw\_split} makes \texttt{HasParamW} applicable and
\texttt{NAut.lab\_loc} identifies $f_v$ with the good-row label.

The first clause of \path{CyclotomicNorms.IsGood} says that
$f_v$ is monic of degree $r$, hence nonzero.  Thus every factor
in (G37.2) is nonzero.  Comparing degrees gives
\[
 r=\sum_i d_i\deg h_{v,i}\le\sum_i d_i a_i=r.
\]
All individual nonnegative deficits vanish, so
$\deg h_{v,i}=a_i$ at every test place.  In particular every
$h_{v,i}$ is nonconstant.  The second \texttt{IsGood} clause
is squarefreeness.  If some $d_i\ge2$, the square of this
nonconstant factor divides $f_v$, a contradiction.  Hence
every $d_i=1$.

If there were at least two blocks, one of their positive
degrees $a_i$ would satisfy $2a_i\le r$.  The factor
$h_{v,i}$ of that fixed degree would divide $f_v$ for every
test place.  The fifth \texttt{IsGood} clause says that for
each $1\le a\le r/2$ some test place has no degree-$a$
factor.  This is a contradiction.  There is therefore one
block, its multiplicity is one, and its degree is $r$.

This proves the exact source theorem
\path{Formalism.MainNotions.SumRepresentations.}
\texttt{parameter\_of\_good\_is\_cuspidal}, in
\path{ExistsInertiaCharpolyFactorOfTameType.lean},
lines 469--588 (the theorem begins at line 522).
The elementary factor argument also appears as
\path{IrreducibleCountingLemmas.Z.goodCuspidal_holds}
in \path{PadicSubsetSums/StableGradedCounts/ZGoodCuspidal.lean}.
It uses the first, second and fifth good-label clauses;
no residual irreducibility theorem is needed for this step.

\proofstep{2. The archimedean obstruction even outside the good block.}
For a scalar character with archimedean parameter $z$, the
literal \texttt{isoInf} definition contributes the multiset
\[
 A_d(z)=\left[\,z+\frac{d-1}{2}-j:\ 0\le j<d\,\right].
                                                               \tag{G37.3}
\]
Let two blocks have positive lengths $d,e$ and the same center
$z$.  If $d,e$ have the same parity, suppose $d\ge e$.
For $0\le j<e$, the integer $j+(d-e)/2$ lies between $0$
and $d-1$ and yields the same entry of $A_d(z)$ as the
$j$th entry of $A_e(z)$.  Thus $A_e(z)$ is contained in
$A_d(z)$ and their multiset sum has a repeated entry.

If $d,e$ have opposite parity, their top entries differ by
$(d-e)/2$, which is not an integer.  Two elements of a single
coset $q+\mathbb Z$ must have integral difference.  Hence
their union is not contained in any one such coset.

The exact source fields \texttt{NBC.infW\_alg} and
\texttt{NBC.embAV\_inf}, in \texttt{DataN.lean}, lines
741--749, together with the last clause of \texttt{IsBCV},
require \texttt{isoInf} both to have no repetitions and
to lie in
$\mathbb Z+(n_V-1)/2$.  The two cases above contradict
these requirements.  A common external twist shifts both
centers by the same number and cannot change either argument.
Changing the ambient rank changes the required coset, but
cannot make a nonintegral difference integral.

A finite-order continuous character of the complex
archimedean multiplicative group is trivial, since that
group is connected.  Therefore a finite-order twist leaves
the scalar archimedean parameter unchanged.  In the literal
source notation the corresponding identity is
\texttt{NAut.infG\_tw} with twisting parameter $0$.
This proves the assertion about a finite $\chi_{269}$ ratio.
The same parity compatibility is also recorded in
\path{SphericalInertData.Pend.part_entries}; no omitted
half-algebraic convention changes (G37.3).

\proofstep{3. A real norm shift is not a free scalar-center parameter.}
This can be checked using the source's finite-place equations,
without adding an unstated axiom about its abstract type
\texttt{CG}.  Suppose two scalar characters are related at
good places by \texttt{lpTw} with real exponent $s$, and
both satisfy \texttt{csd}.  At a conjugate pair of good places
of norm $q>1$, write the first character's eigenvalues as
$\alpha,\alpha^{-1}$.  They are nonzero.  The exact fields
\texttt{NAut.frobEig\_tw} and \texttt{NAut.csd\_unr}
then make the shifted values
$\alpha q^s,\alpha^{-1}q^s$, but also require the second
to be the inverse of the first.  Consequently $q^{2s}=1$,
and $s=0$.  In the usual unitary Hecke-character
interpretation the same conclusion follows from absolute
values.  Real norm twists appearing within the already
specified Speh ladder are accounted for by (G37.3); they
are not an extra independent center shift.

The theorem is restricted to coincident centers and to
real norm shifts preserving these source conditions.
It does not exclude distinct nonparallel algebraic infinity
types.  Those types change the coefficient in the cross
Hom and must be matched to the desired
$\chi_{269}\varepsilon^{-e}$ separately.
\end{proof}

\begin{corollary}[AC-G37-self: the good block cannot hide a cubic self-twist]
\label{AC-G37-self}
Let $F$ be a host field containing $\mathbb Q(\omega)$ and
let $\chi:G_F\to\mu_3$ be nontrivial and unramified at the
two test primes.  Suppose the residual good block satisfies
the explicit conclusion of the source
\path{ArithmeticCylinders.PointsLemma} or
\texttt{TheoremThreePrime}.  Then its restriction to
$G_{F(\ker\chi)}=\ker\chi$ is absolutely irreducible.
At residue characteristic different from $3$, its endomorphism
representation has no $\overline\chi$-line.
For a characteristic-zero integral lift with that residual
property, the endomorphism representation has no $\chi$-line
at any coefficient prime, including $3$.
\end{corollary}
\begin{proof}
The source conclusions are stated in
\path{ArithmeticCylinders.lean}, lines 121--164.
They assert irreducibility on every open normal subgroup
with solvable quotient containing the inertia groups over
one test prime whose bad-prime set does not contain the
coefficient prime.  The two bad-prime sets are disjoint,
so at least one test prime has this property.
The subgroup $H=\ker\chi$ is open and normal, its quotient
is cyclic of order three, and it contains both test-prime
inertia groups by the unramified hypothesis.  The stated
irreducibility conclusion applies.  It is over the algebraic
closure of the residue field, so it is absolute.

Over an algebraically closed field, an endomorphism $T$
transforming by $\chi$ under conjugation commutes with $H$.
An eigenspace of $T$ is then a nonzero $H$-invariant subspace,
so irreducibility forces $T$ to be scalar.  A nonzero scalar
cannot transform by a nontrivial character.  In residue
characteristic different from three, reduction of $\chi$
remains nontrivial, proving the residual assertion.

To pass to characteristic zero, any proper nonzero invariant
subspace over an algebraic closure is defined over a finite
coefficient extension.  Its intersection with the extended
integral lattice is saturated, free of the same rank, and
has a proper nonzero reduction stable under $H$.  This
contradicts absolute residual irreducibility.  The lift
therefore remains absolutely irreducible on $H$, and the
same eigenspace argument applies.  At coefficient three
this argument uses irreducibility of the restriction; it
does not pretend that $\overline\chi$ is nontrivial.
\end{proof}

\begin{lemma}[AC-G37-ann: the source annihilation calculation allows arbitrary ratios]
\label{AC-G37-ann}
Let $R$ be a commutative ring, $C$ a commutative $R$-algebra,
$\epsilon:C\to R$ an $R$-algebra map, and $B$ an associative
$C$-algebra.  Let $r:\Gamma\to B^\times$ and
$a,b:\Gamma\to R^\times$ be homomorphisms.
Suppose $\psi_0,\psi:B\to R$ are ring maps restricting
to $\epsilon$ on $C$ and satisfying
$\psi_0(r_\gamma)=a_\gamma$,
$\psi(r_\gamma)=b_\gamma$.
Suppose $E\in B$, $d\in C$ satisfy
\[
 E^2=dE,\qquad \psi_0(E)=\epsilon(d),\qquad \psi(E)=0,
 \qquad E=\sum_{i=1}^m t_i r_{\gamma_i},\quad t_i\in C.
                                                               \tag{G37.4}
\]
Put $y_\gamma=(a_\gamma^{-1}r_\gamma-1)E$ and let $Y$
be the left $B$-submodule of $B$ they generate.
If $M$ is an $R$-module and $\varphi:Y\to M$ is
$R$-linear with
\[
 \varphi(uy)=\psi(u)\varphi(y),\qquad
 \varphi(y_\gamma)=(b_\gamma a_\gamma^{-1}-1)w
                                                               \tag{G37.5}
\]
for some $w\in M$, then $\epsilon(d)^2w=0$.
There is no cyclotomic or finite-order restriction on $b/a$.
\end{lemma}
\begin{proof}
Set $q=\sum_i t_i a_{\gamma_i}\in C$ and
$z=\sum_i t_i a_{\gamma_i}y_{\gamma_i}\in Y$.
Since $C$ acts centrally on $B$, (G37.4) gives
\[
                 z=E^2-qE=(d-q)E,\qquad Ez=dz.
\]
Using (G37.5) and the two evaluations of $E$,
\[
 \varphi(z)=\sum_i\epsilon(t_i)(b_{\gamma_i}-a_{\gamma_i})w
           =(\psi(E)-\psi_0(E))w=-\epsilon(d)w.
\]
On the other hand $\varphi(Ez)=\psi(E)\varphi(z)=0$,
whereas $\varphi(dz)=\epsilon(d)\varphi(z)
=-\epsilon(d)^2w$.  The equality $Ez=dz$ proves the result.

This is the character-ratio part of the source theorem
\path{TypePrimeCongruences.CongruenceModules.}
\texttt{congruence\_annihilates\_character\_part}, in
\path{CongruenceAnnihilatesCharacterPart.lean}, lines
368--426.  Its source elements are $E=\texttt{Eel}$,
$d=\texttt{del51}$ and the evaluations
$\texttt{psiChi}$ at the first and selected character.
The proof above isolates exactly the identities used;
it invokes no \texttt{NData.chi\_def}.
It therefore remains valid for a mixed finite-character
ratio once those identities have actually been constructed.
It does not construct the needed idempotent-choice data,
local extension conditions, or global Selmer comparison.
\end{proof}

\paragraph{The smaller remaining comparison.}
Two separate regular packets related by a finite cubic
twist preserve their individual infinity types and avoid
the simultaneous-ladder obstruction.  A comparison of
their integral Hecke lattices is a different problem from
putting both ladders in one \texttt{IsBCV} parameter.
The existing \texttt{chi\_def} and \texttt{sha\_eq}
still give the pure-Tate length input described in AC-G29.
The preceding annihilation lemma removes one algebraic
substep from a proposed relative comparison, but does not
identify its indices with \texttt{LinkN} or infer any
new Selmer upper bound from an index upper bound.

\paragraph{Geometric dictionary.}
The passage from the finite cubic cover on the open where it
is unramified to its character and kernel uses the already
audited Galois-category dictionary of SGA~1, Expos\'e~V,
\S\S4 and~7.  Restriction to that kernel is base change to
the corresponding cover.  The centered-weight obstruction
above is the explicit multiset calculation in the source's
archimedean realization; no additional SGA~6 comparison
theorem or unverified archive locator is invoked.
